\documentclass[11pt]{article}
\usepackage[a4paper,margin=25mm]{geometry}
\usepackage[T1]{fontenc}
\usepackage{lmodern,amsmath,amssymb,amsthm,mathtools,microtype}
\usepackage[hidelinks]{hyperref}
\usepackage{enumitem}
\setlist{nosep,leftmargin=*}
\newtheorem{theorem}{Theorem}
\newtheorem{lemma}[theorem]{Lemma}
\newtheorem{corollary}[theorem]{Corollary}
\theoremstyle{remark}\newtheorem{remark}[theorem]{Remark}
\newcommand{\Pp}{\mathbb P}
\newcommand{\E}{\mathbb E}
\newcommand{\N}{\mathbb N}
\newcommand{\supp}{\operatorname{supp}}
\newcommand{\one}{\mathbf 1}
\newcommand{\fall}[2]{(#1)_{\underline{#2}}}
\title{Subsequence complete condensation\protect\\for zero-radius allocations}
\author{}
\date{17 September 2026}
\begin{document}
\maketitle
\vspace{-1.5em}
\begin{center}\small Complete proof candidate; independent verification pending.\end{center}
\begin{abstract}
For a fixed nonnegative weight sequence with positive constant term and zero
radius of convergence, consider weighted allocations of $m_n$ balls into $n$
boxes, where $m_n/n\to\lambda\in(0,\infty)$. We give an argument that every
admissible full path has a subsequence on which the largest box contains
$(1-o_{\Pp}(1))m_n$ balls. No bounded-increment assumption on $m_n$ is used.
The two main ingredients are an occupation-count insertion inequality that
retains the entire small-box partition factor, and a geometric-record
summation in which the packing penalties are charged once per record
plateau. The asserted conclusion answers Janson's Problem 18.40 negatively
and is stronger than the failure of full-sequence sublinear condensation.
\end{abstract}

\section{Model and statement}
Let $w=(w_j)_{j\ge0}$ be fixed, with $w_0>0$, $0\le w_j<\infty$, and
\begin{equation}\label{eq:radius}
 \rho=0,\qquad\text{equivalently}\qquad
 \limsup_{j\to\infty}w_j^{1/j}=\infty.
\end{equation}
For nonnegative integers $m,n$, set
\[
 Z(m,n)=\sum_{y_1+\cdots+y_n=m}\prod_{v=1}^n w_{y_v}.
\]
A pair is \emph{admissible} when $Z(m,n)>0$. On an admissible pair define
\[
 \Pp(Y=y)=\frac{1}{Z(m,n)}\prod_{v=1}^n w_{y_v},
 \qquad Y_{(1)}\ge Y_{(2)}\ge\cdots
\]
for the ordered occupancies. The weights do not depend on $n$.

\begin{theorem}\label{thm:main}
Suppose \eqref{eq:radius} holds and $m_n$ is an integer sequence such that
$(m_n,n)$ is admissible for every sufficiently large integer $n$ and
$m_n/n\to\lambda\in(0,\infty)$. There is an increasing sequence $n_r$ such that
\begin{equation}\label{eq:conclusion}
 \frac{Y_{(1)}(m_{n_r},n_r)}{m_{n_r}}
 \xrightarrow{\Pp}1.
\end{equation}
The same conclusion holds when the indices range over a set
$\mathcal N=\{N_1<N_2<\cdots\}$ with $N_{j+1}/N_j\to1$, with admissibility
and density assumed on that set.
\end{theorem}

\begin{corollary}\label{cor:consequences}
Under the theorem's assumptions,
\[
 \limsup_{n\in\mathcal N}\E\frac{Y_{(1)}(m_n,n)}n=\lambda.
\]
Along a subsequence, $Y_{(1)}/n\to\lambda$ in probability and in $L^1$,
$(m_n-Y_{(1)})/n\to0$ in probability and in $L^1$, and $Y_{(2)}/n\to0$ in
probability. In particular $Y_{(1)}/n$ cannot converge to zero along the full
path.
\end{corollary}

Janson's Problem 18.40 asks whether full-sequence convergence to zero is
possible in the regime $\lambda>\nu=0$ \cite[Section 18.6, Problem 18.40]{Janson}.
For a nontrivial weight sequence with $w_0>0$, this is the zero-radius regime.
The theorem permits arbitrary jumps and oscillations of $m_n$, subject only
to the stated limiting density. A preselected lacunary subsequence is not a
full path; the source's Example 18.38 already shows that such subsequences
can have a sublinear maximum.

Multiplying every weight by $1/w_0$ leaves the allocation law unchanged, so
we normalize $w_0=1$. Write $d=\gcd\{j\ge1:w_j>0\}$. Dividing every supported
size and every admissible total by $d$ preserves the ratio $Y_{(1)}/m$ and
zero radius and replaces $\lambda$ by $\lambda/d$. It therefore suffices to
prove the theorem for span one. We make this reduction throughout the proof.

\section{Insertion without losing the small-box factor}
For an integer $u\ge0$ write
\[
 Z_u(t,n)=\sum_{\substack{x_1+\cdots+x_n=t\\0\le x_v\le u}}
             \prod_{v=1}^n w_{x_v}.
\]
The following finite inequality is the reason no control of the jumps of
$m_n$ is needed.

\begin{lemma}[Occupation-count insertion]\label{lem:insertion}
Let $u<k$, $w_k>0$, and $m=t+k+r<2k$. Suppose $r$ has a fixed representation
by $R$ supported positive sizes at least $K$, together with $C$ other
supported positive sizes. All these sizes are less than $k$. Let $W_r$ be
the product of their weights. If
\begin{equation}\label{eq:insertionconditions}
 nK\ge2m,\qquad t+1+R+C\le n/2,
\end{equation}
then
\begin{equation}\label{eq:insertion}
 Z(m,n)\ge\frac{n}{2^{C+1}}\,w_k W_r Z_u(t,n).
\end{equation}
Repeated sizes, including sizes occurring in both parts of the
representation, are allowed.
\end{lemma}

\begin{proof}
Describe a background configuration by its occupation counts $a_j$, where
$\sum a_j=n$ and $\sum j a_j=t$, and $a_j=0$ for $j>u$. The total weight of
that count vector, including all coordinate permutations, is
\[
 \mathcal W(a)=\frac{n!}{\prod_j a_j!}\prod_j w_j^{a_j}.
\]
Set $Q=1+R+C$. Use $Q$ empty coordinates to insert one size-$k$ box and the
chosen representation of $r$. If $r_j$ is the multiplicity of size $j$ in
that representation, the resulting count vector is
\[
 b_0=a_0-Q,\quad b_k=1,\quad b_j=a_j+r_j\quad(0<j<k).
\]
The map $a\mapsto b$ is injective. Since $a_0\ge n-t$, its weight ratio is
\begin{equation}\label{eq:typeratio}
 \frac{\mathcal W(b)}{\mathcal W(a)}
 =w_kW_r\frac{\fall{a_0}{Q}}
 {\displaystyle\prod_{0<j<k}\prod_{\ell=1}^{r_j}(a_j+\ell)}.
\end{equation}
The numerator is at least $(n/2)^Q$ by \eqref{eq:insertionconditions}.
Each denominator factor assigned to one of the $R$ large pieces is at most
$m/K$, since the final configuration has total $m$. Each of the other $C$
factors is at most $n$. When the two classes share a size, assign their
factors in either fixed order; both bounds still apply. Thus the ratio in
\eqref{eq:typeratio} is at least
\[
 w_kW_r\frac{n}{2^{C+1}}
       \left(\frac{nK}{2m}\right)^R
 \ge w_kW_r\frac{n}{2^{C+1}}.
\]
Summing over the distinct background count vectors gives
\eqref{eq:insertion}. Weights of sizes not used in the construction only
increase the original denominator.
\end{proof}

\section{Geometric records and two summation bounds}
Fix a real base $a>1$ and a constant $b>2$. For every sufficiently large
integer $i$ define
\begin{equation}\label{eq:records}
 B_i=\max_{\substack{1\le s\le\lfloor a^i\rfloor\\w_s>0}}
       \frac{\log w_s+b\log s}{s},\qquad
 s_i=\min\operatorname*{arg\,max}_{s}\frac{\log w_s+b\log s}{s},
 \qquad p_i=\frac{b\log s_i}{s_i}.
\end{equation}
Then $B_i$ is nondecreasing and tends to infinity. Also $s_i\to\infty$.
Set $d_i=B_i-B_{i-1}$ and call an index \emph{active} if $d_i>0$. At an
active index $J$,
\begin{equation}\label{eq:active}
 a^{J-1}<s_J\le a^J,\qquad
 \log w_{s_J}=s_JB_J-b\log s_J.
\end{equation}
At inactive indices the smallest maximizing index remains unchanged.
Consequently $s_i$ and $p_i$ are constant on each plateau of $B_i$.

\begin{lemma}[Summable plateau charges]\label{lem:plateau}
Fix a sufficiently large starting index $i_0$ and an integer $D\ge1$.
If $\mathcal A$ is the set of active indices after $i_0$, then
\begin{equation}\label{eq:plateausum}
 \sum_{J\in\mathcal A}\ \sum_{i=i_0}^{J-D} a^{i-J}p_i<\infty.
\end{equation}
More precisely, this sum is at most
\[
 \frac{a}{(a-1)^2}
 \left(p_{i_0}+\sum_{\ell\in\mathcal A}p_\ell\right),
\]
and the parenthesized sum is finite.
\end{lemma}

\begin{proof}
For an active $\ell$, \eqref{eq:active} gives
\[
 p_\ell\le ba\ell(\log a)a^{-\ell},
\]
which is summable. Consider a plateau starting at $\ell$ and ending just
before the next active index $v$. A term from this plateau can contribute
to \eqref{eq:plateausum} only at an active $J\ge v$: the starting active
index $J=\ell$ has cutoff $J-D<\ell$. Its total contribution is at most
\[
 p_\ell
 \left(\sum_{i=\ell}^{v-1}a^i\right)
 \left(\sum_{J=v}^{\infty}a^{-J}\right)
 \le \frac{a}{(a-1)^2}p_\ell.
\]
The same bound applies to the initial plateau at $i_0$. There are always
later active indices because $B_i\to\infty$. Summing proves the lemma.
\end{proof}

\begin{lemma}[Contractive memory kernel]\label{lem:kernel}
For $D\ge2$ put
\begin{equation}\label{eq:theta}
 C_a=2a^2,\qquad
 \Theta(a,D)=2a^{3-D}
 \left(\frac{D-1}{a-1}+\frac{1}{(a-1)^2}\right).
\end{equation}
For every nondecreasing sequence $B_i$ and every finite upper limit $N$,
\begin{equation}\label{eq:kernelsum}
 C_a\sum_{\substack{i_0<J\le N\\J\ {\rm active}}}
       \sum_{i=i_0}^{J-D}a^{i-J}(B_{J-1}-B_i)
 \le\Theta(a,D)\sum_{\ell=i_0+1}^{N}d_\ell.
\end{equation}
In particular $D$ can be chosen so that $\Theta(a,D)<1/4$.
\end{lemma}

\begin{proof}
Expand $B_{J-1}-B_i=\sum_{\ell=i+1}^{J-1}d_\ell$. All terms are
nonnegative, so extend the sum over $J$ to all larger integers and the
lower limit on $i$ to $-\infty$ to bound each coefficient. For a fixed
$\ell$ and lag $t=J-\ell\ge1$, that coefficient, before multiplication by
$C_a$, is
\[
 \sum_{i\le\min(J-D,\ell-1)}a^{i-J}
 =\begin{cases}
  a^{1-D}/(a-1),&1\le t\le D-1,\\
  a^{-t}/(a-1),&t\ge D.
 \end{cases}
\]
Summing over $t$ gives
$a^{1-D}((D-1)/(a-1)+1/(a-1)^2)$. Multiplying by $2a^2$ gives
\eqref{eq:theta}. Exponential decay in $D$ proves the last assertion.
\end{proof}

\section{A supported representation of every small residual}
The span-one reduction provides a finite set of positive supported sizes
whose gcd is one. Its nonnegative additive semigroup contains every integer
above some $h_*\ge1$. Here is a direct justification. Choose one member
$s$ of the finite set. The nonnegative residues generated modulo $s$ form
a subgroup of the finite cyclic group and hence all residues. Choose one
representable integer in each residue class; adding copies of $s$ then
represents every sufficiently large integer.

Fix an integer $K\ge2$. Enlarge $i_0$ so that, for all $i\ge i_0$,
\begin{equation}\label{eq:start}
 s_i\ge K,\qquad B_i\ge0,\qquad a^{i_0}\ge h_*.
\end{equation}
Choose once and for all a supported representation for each integer in
\[
 h_*\le v<2a^{i_0}.
\]
Let $C_*$ be a common bound on their numbers of positive parts, let
$M=\lceil2a^{i_0}\rceil$, and choose $C_0\ge0$ such that the logarithm of
each of their product weights is at least $-C_0$. These are finite constants.

\begin{lemma}[Greedy residual bound]\label{lem:greedy}
Suppose $r\ge h_*$ and
\[
 I(r):=\left\lfloor\log_a(r/2)\right\rfloor\le J-D,
 \qquad D\ge2.
\]
There is a representation of $r$ using supported parts at least $K$, and
at most $C_*$ additional supported parts, with product weight $W_r$ satisfying
\begin{equation}\label{eq:greedy}
 B_{J-1}r-\log W_r
 \le 2\sum_{i=i_0}^{J-D}a^{i+1}
       (B_{J-1}-B_i+p_i)+MB_{J-1}+C_0.
\end{equation}
All parts are at most $r$. The construction uses at most one stage at each
geometric index $i$.
\end{lemma}

\begin{proof}
With current residual $x$, stop if $x<2a^{i_0}$ and use its fixed terminal
representation. Otherwise set
\[
 i=\lfloor\log_a(x/2)\rfloor,\quad
 q=\left\lfloor\frac{x-h_*}{s_i}\right\rfloor,\quad
 x'=x-qs_i.
\]
Here $i\ge i_0$, $q\ge1$, and
\[
 h_*\le x'<h_*+s_i\le2a^i,\qquad qs_i\le x<2a^{i+1}.
\]
Thus the next geometric index is strictly smaller than $i$. Add $q$ copies
of $s_i$ and repeat. The procedure terminates in the fixed terminal range.
At stage $i$ the logarithmic weight per unit mass is $B_i-p_i$.
Since $B_{J-1}\ge B_i$ for every index used, that stage's deficit against
$B_{J-1}$ is at most
$2a^{i+1}(B_{J-1}-B_i+p_i)$. The terminal deficit is at most
$MB_{J-1}+C_0$. Unused summands are nonnegative, giving \eqref{eq:greedy}.
\end{proof}

\section{Proof for arbitrary total-mass paths}
We first prove the full-integer-index version of Theorem~\ref{thm:main}.
Fix a desired error $\varepsilon\in(0,1)$ and choose $a>1$ sufficiently
close to one. Fix $b>2$, say $b=4$. Choose $D\ge2$ with
$\Theta(a,D)<1/4$. Then choose a fixed $\delta>0$ sufficiently small that
\begin{equation}\label{eq:parameters}
 \delta<\min\left\{\frac{a^{-D}}{100},\frac1{100(1+\lambda)},\frac1{100}\right\},
 \qquad \frac1{a(1+4\delta)}>1-\varepsilon.
\end{equation}
For example first choose $a<1/(1-\varepsilon)$, and then decrease $\delta$
if necessary. All these parameters are fixed before taking record indices
to infinity.

Take an integer $K\ge4(1+\lambda)$ and choose the constants $i_0,h_*,C_*,M,C_0$
as in Section~4. At each sufficiently large active $J$ set
\begin{equation}\label{eq:choice}
 k=s_J,\quad u=\lfloor a^{J-1}\rfloor,\quad h=B_{J-1},\quad
 n=n(J)=\left\lceil\frac{(1+4\delta)k}{\lambda}\right\rceil,
 \quad m=m_{n(J)}.
\end{equation}
The prescribed path, regardless of its jumps, satisfies
\begin{equation}\label{eq:path}
 \frac nk\longrightarrow\frac{1+4\delta}{\lambda},\qquad
 \frac mk\longrightarrow1+4\delta,
 \qquad 3\delta k\le m-k\le5\delta k
\end{equation}
for all sufficiently large active $J$. In particular $m<2k$.

\subsection{A typical background and exact cancellation}
Define an auxiliary distribution supported on $\{0,\ldots,u\}$ by
\[
 q_j=w_je^{-hj},\quad A=\sum_{j=0}^{u}q_j,\quad p(j)=q_j/A.
\]
From \eqref{eq:records},
\[
 q_0=1,\qquad q_j\le j^{-b}\quad(1\le j\le u).
\]
As $h\to\infty$, dominated convergence gives
\begin{equation}\label{eq:auxmean}
 A\to1,\qquad \mu=\sum_{j=0}^{u}j p(j)\to0.
\end{equation}
The dominating series for the mean is $\sum j^{1-b}<\infty$.
For $S_n$ the sum of $n$ independent copies of $p$, Markov's inequality and
\eqref{eq:path} show $\Pp(S_n\le\delta k)\to1$. Hence there is an integer
$t\in[0,\lfloor\delta k\rfloor]$ with
\begin{equation}\label{eq:atom}
 \Pp(S_n=t)\ge\frac1{2(\lfloor\delta k\rfloor+1)}.
\end{equation}
Let $r=m-k-t$. Then $2\delta k\le r\le5\delta k$. In particular $r\ge h_*$
eventually, and \eqref{eq:parameters} gives $I(r)\le J-D$.
Use the representation from Lemma~\ref{lem:greedy}.

All its pieces are less than $k$. If it has $R$ large pieces and $C\le C_*$
terminal pieces, then $R+C\le r$, and
\[
 t+1+R+C\le6\delta k+1<n/2
\]
for all sufficiently large $J$. Also $nK\ge2m$ eventually. Thus
Lemma~\ref{lem:insertion} applies. With
\[
 P_J:=\Pp_w(Y_{(1)}(m,n)\le u)=\frac{Z_u(m,n)}{Z(m,n)},
\]
it yields
\begin{align}
 P_J
 &\le\frac{2^{C_*+1}}{n}
       \frac{Z_u(m,n)}{w_kW_r Z_u(t,n)}\notag\\
 &\le\frac{2^{C_*+2}(\lfloor\delta k\rfloor+1)}{n}
       \exp\{hk-\log w_k+hr-\log W_r\}.\label{eq:probabilitybound}
\end{align}
Indeed,
$Z_u(t,n)=e^{ht}A^n\Pp(S_n=t)$ and $Z_u(m,n)\le e^{hm}A^n$.
The same $A^n$ occurs in both terms and cancels \emph{exactly}. There is no
loss of the form $A^L$ for padded reservoir coordinates. The prefactor in
\eqref{eq:probabilitybound} is bounded by a fixed constant $C_1$.

\subsection{Assuming a positive lower bound leads to a contraction}
Suppose that, for some $\eta>0$, $P_J\ge\eta$ at every sufficiently large
active $J$. By \eqref{eq:active}, \eqref{eq:greedy}, and
\eqref{eq:probabilitybound},
\begin{align}
 d_J\le{}&C_a\sum_{i=i_0}^{J-D}a^{i-J}
                 (B_{J-1}-B_i+p_i)\notag\\
 &+Ma^{1-J}B_{J-1}
   +a^{1-J}\bigl(bJ\log a+C_2\bigr),\label{eq:recurrence}
\end{align}
where $C_2$ is a finite constant allowed to depend on $\eta$.
We used $k>a^{J-1}$ and $\log k\le J\log a$. Enlarge $C_2$ to be
nonnegative when necessary.

Choose $J_0$ beyond all preceding thresholds and so large that
\begin{equation}\label{eq:kappa}
 \kappa:=\frac{Ma^{1-J_0}}{a-1}<\frac14.
\end{equation}
For a finite upper limit $N$ write
$S_N=B_N-B_{J_0-1}=\sum_{J=J_0}^N d_J$. Sum \eqref{eq:recurrence} over
active $J\in[J_0,N]$. Inactive increments are zero. Lemma~\ref{lem:kernel}
bounds the memory term by $\Theta(a,D)S_N$ plus a fixed contribution from
the earlier increments. Lemma~\ref{lem:plateau} bounds all $p_i$ terms by
a constant independent of $N$. The final term in \eqref{eq:recurrence} is
summable, since $\sum_J (J+1)a^{-J}<\infty$.

For completeness, the remaining terminal term satisfies
\begin{align*}
 M\sum_{J=J_0}^N a^{1-J}B_{J-1}
 &\le\frac{MB_{J_0-1}a^{2-J_0}}{a-1}
    +M\sum_{\ell=J_0}^{N-1}d_\ell
                   \sum_{J=\ell+1}^{\infty}a^{1-J}\\
 &\le C_3+\frac{Ma^{1-J_0}}{a-1}S_N
 =C_3+\kappa S_N.
\end{align*}
Therefore a constant $C_4<\infty$, independent of $N$, satisfies
\begin{equation}\label{eq:contradiction}
 S_N\le(\Theta(a,D)+\kappa)S_N+C_4.
\end{equation}
Both coefficients are less than $1/4$, so $S_N$ is bounded. This contradicts
$B_N\to\infty$. We have proved
\begin{equation}\label{eq:liminf}
 \liminf_{\substack{J\to\infty\\J\ {\rm active}}}P_J=0.
\end{equation}

\subsection{Near-total mass and the order of quantifiers}
By \eqref{eq:active} and \eqref{eq:path},
\[
 \liminf_{\substack{J\to\infty\\J\ {\rm active}}}\frac{u}{m}
 \ge\frac1{a(1+4\delta)}>1-\varepsilon.
\]
Thus \eqref{eq:liminf} implies that for every $\eta>0$ there are arbitrarily
large $n$ with
\[
 \Pp_w\bigl(Y_{(1)}(m_n,n)\le(1-\varepsilon)m_n\bigr)<\eta.
\]
Apply this assertion separately with $\varepsilon=1/r$ and $\eta=1/r$,
choosing the sizes strictly increasingly. This proves \eqref{eq:conclusion}.
The weights and the path remain fixed; only the auxiliary parameters depend
on the desired error before each limiting argument.

If the admissible indices are $\mathcal N=\{N_1<N_2<\cdots\}$ with
$N_{j+1}/N_j\to1$, replace the ceiling in \eqref{eq:choice} by the first
$N_j\ge(1+4\delta)k/\lambda$. The first two limits in \eqref{eq:path} are
unchanged. Every subsequent step is identical. This proves the second
assertion of Theorem~\ref{thm:main}, and the span reduction at the start
restores arbitrary supports.

Finally $0\le Y_{(1)}/m_n\le1$, so convergence in probability along the
chosen subsequence implies convergence in $L^1$. Multiplication by
$m_n/n\to\lambda$, and $Y_{(2)}\le m_n-Y_{(1)}$, prove
Corollary~\ref{cor:consequences}. This completes the argument.

\section{Scope and checks}
The conclusion is \emph{subsequence} complete condensation for every fixed
full mass path, not condensation at every size and not
$\Pp(Y_{(1)}=m_n)\to1$. The distinction is necessary. For the source's sparse
factorial support, the maximum is deterministically sublinear at the
lacunary sizes of Example 18.38 \cite{Janson}. For full factorial weights,
small residual branches may remain even when their total is negligible
relative to the giant \cite[Theorem 2.4]{JJS}.

For simply generated plane trees, the maximum outdegree has the same law
as the maximum allocation with $m=n-1$: the cyclic lemma groups every
allocation sequence of total $n-1$ into an orbit of $n$ rotations containing
exactly one valid tree sequence, while preserving product weights and the
maximum. The admissible large tree sizes lie on a fixed arithmetic lattice,
so the dense-index extension applies. The preceding bounded-jump proof is
therefore subsumed rather than contradicted.

The positive-boundary-mean regime of Problem 18.39 is not addressed here.
The use of zero radius is precise: it makes $B_i\to\infty$ and makes the
auxiliary background mean tend to zero. No unproved large-deviation,
local-limit, monotonicity, or regular-variation result is invoked.

The accompanying \texttt{check\_full\_path.py} uses exact integer and rational
arithmetic to test the occupation-count ratio, insertion lower bounds,
the geometric kernel constant, plateau charging including long plateaus,
and finite greedy representations and energy bounds. These checks do not
constitute a formalization or an independent audit of the asymptotic proof.
Publication priority has not been established by an exhaustive literature
review.

\begin{thebibliography}{9}
\bibitem{Janson}
S. Janson,
\emph{Simply generated trees, conditioned Galton--Watson trees, random
allocations and condensation},
Probability Surveys \textbf{9} (2012), 103--252.
DOI: 10.1214/11-PS188. arXiv:1112.0510.

\bibitem{JJS}
S. Janson, T. Jonsson and S. \"O. Stef\'ansson,
\emph{Random trees with superexponential branching weights},
Journal of Physics A: Mathematical and Theoretical \textbf{44} (2011), 485002.
arXiv:1104.2810.
\end{thebibliography}
\end{document}
